\begin{textsample}{NEG Dim 3 – ce – Score -3.00 – ce/ce\_em\_61.txt}  \label{ex:f3_neg_215}
9.5.3 Itinerário para a área de Ciências da Natureza e suas Tecnologias Na perspectiva do novo ensino \textbf{médio}, por conta da necessidade de um maior envolvimento das/os alunas/os no processo de ensino e aprendizagem, é buscado nestes, o desenvolvimento de \textbf{competências}, habilidades, atitudes e valores para a melhor formação do cidadão do futuro.
Sendo assim, as Diretrizes Curriculares Nacionais para o Ensino Médio ( DCNEM ), atualizadas pelo Conselho Nacional de Educação ( CNE ) em novembro de 2018, indicam que os currículos dessa etapa de ensino devem ser compostos por uma formação geral básica e itinerários formativos.
Neste contexto, os itinerários formativos representam, um conjunto de situações e atividades educativas que as/os alunas/os podem escolher conforme seu interesse, para aprofundar e ampliar aprendizagens em uma ou mais Áreas de Conhecimento e/ou na Formação Técnica e Profissional.
Para tanto, os itinerários formativos ficam compostos por eixos estruturantes, que buscam integrar os diferentes conhecimentos associados a sua realidade para oportunizar aos alunas/os, experiências que promovam sua formação pessoal, profissional e cidadã. 9.5.3.1 Perfil do Egresso Após a participação efetiva nos itinerários, as/os alunas/os deverão ter desenvolvido habilidades e \textbf{competências} que os tornem capazes de realizarem as devidas conexões entre os saberes científicos, culturais, sociais e econômicos trabalhados nesta área do conhecimento, bem como potencializar a sua capacidade de intervenção social para o exercício da cidadania, transformando para melhor, a sua vida e da comunidade. 9.5.3.2 Objetivos do itinerário Esse itinerário tem como objetivo principal a ampliação dos conhecimentos trabalhados na Formação Geral Básica ou seu aprofundamento, dinamizar a interação do saber entre as áreas do conhecimento e ainda, despertar as conexões necessárias ao desenvolvimento intelectual, moral e social dos/as discentes, de modo a serem capazes de interferir construtivamente no contexto social em que visem a melhorar sua qualidade de vida, da família e da comunidade. 9.5.3.3 Detalhamento do itinerário formativo As unidades curriculares sugeridas neste itinerário formativo pode ser trabalhadas de acordo com a necessidade, condições e proposta pedagógica da escola, focando atender conjuntura local, propiciando o reconhecimento da natureza e as possibilidades desta e subsidiando o desenvolvimento de atividades que contribuam para o crescimento individual e coletivo.
O caminho a seguir, deve ser possível de adaptação à realidade escolar, podendo haver deslocamentos entre os semestres em todo o Ensino Médio.
Estas propostas foram pensadas para serem distintas e harmônicas, dentro do contexto que a comunidade necessita, atendendo suas especificidades e contemplando a diversidade de conhecimentos possíveis de serem organizados e trabalhados pelo menos na área do conhecimento, valorizando a contextualização e interdisciplinaridade.
As/Os docentes têm a liberdade de livre escolha entre as eletivas disponíveis pois, estas se apresentam inter relacionadas com a trilha do itinerário formativo, as unidades curriculares e a várias atividades e projetos propostos, tornando o conhecimento mais prático e efetivo.
No que se refere às Eletivas, as unidades propostas estão em sua maioria no Catálogo de Componentes Eletivos - Ano 2021.
Elas são de livre escolha dos estudantes, sendo reflexo de sua necessidade de aprofundamento de conhecimentos relacionados à Área de Natureza e/ou de diversificação de sua formação.
Para a 1a série, a carga horária das eletivas será : no primeiro semestre 10h/a e no segundo 4 h/a, cada eletiva possui carga horária de 2 h/a semanais.
A trilha sugerida para o segundo semestre da 1a série é uma experimentação, ou seja, objetiva que o estudante conheça a área de conhecimento e caso tenha seu interesse despertado, possa seguir consolidando seus conhecimentos nas trilhas de aprofundamento subsequentes.

% matched lemmas: competência, médio
\end{textsample}
